\section{演化策略：赛博寄生与复杂度棘轮}

智能要改变的现实，在它到达之前就已经建成。道路、楼宇、医院和工厂各有既存结构，也嵌入法规、保险、习惯与责任安排。模型要通过这些接口工作，部署成本便取决于需要改动多大的范围。换一个部件是维修，换一套承重结构是重建。单点能力足够以后，还要比较适配旧结构与重建环境的总成本。

\subsection{覆盖层策略与重构的三重代价}

\subhead{重构的三重代价}

理论上，最干净的道路是重构底层。为自动驾驶重新设计道路、车辆、信号、保险和交通规则；为机器人重新设计楼宇、工具、仓库、家庭和工厂；为智能医疗重新设计医院流程、病历结构、责任制度和支付体系。这样，物理世界本身会变得更适合机器感知、机器行动和机器决策，系统会更低熵，接口会更清楚，反馈会更稳定。但这条路太贵。

彻底重构原子世界需要克服三重代价。第一是资源势垒：道路、工厂、医院、港口、电网、住宅和城市管网已经沉淀了巨大的材料、资本、空间和劳动，推倒重建，不是修改代码，而是废弃既有物质结构并重新支付能量账单。第二是时间刚性：基础设施有自己的寿命周期，施工、审批、迁移、协调和行为改变都不能像软件版本一样即时更新。第三是拓扑耦合：道路不只是道路，它连接保险、法规、物流、地产、职业分工和居民生活，医院不只是建筑，它连接医生责任、患者信任、医保支付和法律边界。改动底层，就会牵动整个系统。

水泥、钢铁与设备的大规模年度产出，逐年沉积为建筑、道路、厂房和网络。USGS与世界钢铁协会的2026年报告仍显示，全球水泥与粗钢产出分别处在数十亿吨与近二十亿吨的量级。\cite{usgs2026cement,worldsteel2026figures} 存量改造取决于寿命、替换率、建设能力和相邻系统；不能把某一年的产量直接等同于可同时改造的存量。

这三重代价不是新出现的东西，它们正是上一章那三项在建成环境中的具体形式。资源势垒对应已经支付出去的材料、能量和空间成本，时间刚性对应基础设施无法压缩的过程时长，拓扑耦合对应道路、法规、保险和职业分工之间已经形成的连接密度。一个已经建成的东西，它的黏滞不是被谁后来加上去的，而是在建成那一刻就写进了结构。因此，判断一项智能改造是否可行，不能只看它在符号空间里节省了多少，还要看它在多大范围内要求这套旧结构重新支付一次。

\subhead{覆盖层策略}

所以文明通常不会优先选择重构，而会选择覆盖。所谓\term{覆盖层策略}，是指在不彻底改变旧物理骨架的情况下，在其表面叠加一层新的比特控制层。道路不变，加导航、摄像头、雷达、高清地图和云端训练；工厂不变，加传感器、预测维护、数字孪生和调度平台；医院不变，加影像识别、病历摘要、智能分诊和审核系统；城市不变，加摄像头、平台、算法、数据接口和“智慧治理”。原子骨架保持沉重，数字薄膜覆盖其上。

这是一种赛博寄生。这里的“寄生”不是道德意义上的贬义，而是一个结构描述，也是一种演化策略：新的智能层并不独立存在，也不立即替换旧层，而是依附在旧物理系统之上，利用其已经形成的空间、能源、设备、制度和人员网络，获得进入现实的行动位置，并试图改变旧系统的行为。比特层像一层神经组织，覆盖在旧世界的骨架上；它不替换骨骼，却试图让骨骼行动得更快、更准、更可调度。

这种策略之所以普遍，是因为它符合最小代价原则。彻底重构世界太重，局部覆盖则轻得多：加装传感器比重建道路便宜，开发调度系统比重建工厂便宜，训练识别模型比重写医疗制度便宜，给机器人增加感知和控制能力比重新设计全部生活空间便宜。覆盖层是一种绕过高黏滞底层的次优解，它不解决底层沉重性，而是在沉重性上方寻找可操作空间。短期看，覆盖层非常有效：用户端更顺滑，管理端更可视，调度更实时，风险更可监控，局部效率确实提高。一个旧系统被数字化后，看起来像是变轻了，司机有导航，工厂有预测维护，医生有辅助诊断，城市有监控平台，仓库有自动分拣，许多原本依赖人工经验和局部判断的流程，被新的比特层重新组织。

覆盖并不是数字时代才有的策略。19 世纪末到 20 世纪初，工厂的动力从蒸汽机集中传动改为电动机分散驱动，技术替换在 1890 年代就已大规模展开，但生产率的大幅提升要等到厂房布局、传动方式、机器排布和生产组织都围绕分散动力重排之后才出现，中间隔了二三十年。集装箱的推广同样如此：1956 年第一次集装箱运输之后，码头、起重机、卡车底盘、铁路平车、仓储布局、工会规则和关税制度都要跟着改，港口装卸区被整体搬迁，直到 1970 年代集装箱化才真正铺开。两者都说明，新层接入旧骨架时，瓶颈往往不在新层本身，而在旧骨架愿不愿意为它改变。

近年的例子更容易看清覆盖层的收益边界。2019 年起全国推广高速公路电子不停车收费，江苏、山东、四川、重庆的试点统计显示，客车通过省界的时间从 15 秒降到 2 秒，货车从 29 秒降到 3 秒；2020 年 1 月 1 日，全国 29 个联网省份的 487 个省界收费站全部取消。收费站这个环节确实被比特层改写了，但车道数、匝道容量、主线通行能力和服务区容量都没有随之改变。节假日的高峰流量仍然要在同样的物理通道里消化，拥堵从收费站前移到了主线和出入口。当收费站不再是瓶颈，瓶颈会回到骨架本身上来。这笔收益是真的，但它只在骨架允许的范围内兑现。

旧世界之所以沉，还因为它把一批历史决定固化成了物理尺寸。铁路轨距在 19 世纪被定在 1435 毫米，此后隧道、桥梁、站台和车辆全部按这个宽度建造。集装箱在 1960 年代被定在 20 英尺和 40 英尺两个长度，港口的起重机跨距、堆场分区、卡车底盘和铁路平车随之定型。楼梯踏步的高度上限、住宅门洞的宽度档位、插座形状、交流电频率、键盘字母排布，都是按当时的人体尺度、工艺条件和行业妥协定下来的。这些东西今天还在用，不是因为它们最优，而是因为改动它们的代价要按原子世界的方式支付，并且要由整条链条分摊。

覆盖层策略之所以普遍，还因为它把成本记在了另一个账本上。加装传感器、部署摄像头、上线调度平台，这些支出落在当年的资本开支和运维预算里，看得见也算得清；而旧骨架没有为机器感知和机器行动改变过这件事，不会被计入任何一项成本。前者是显性支出，后者是隐性欠账。短期看，覆盖层的收益立刻显现，一个旧系统被数字化之后确实变顺滑了。代价要等到覆盖层之间开始互相依赖时才出现，而那时它已经变成了维持运行的必要条件。

\subsection{复杂度棘轮、兼容税与落地税}

\subhead{复杂度棘轮}

覆盖层需要传感、数据、通信、权限、监控、维护与异常处理。新增接口若同时要求旧结构作出配套，局部收益可能带来新的耦合；结构又依赖更多结构来维持时，就形成\term{复杂度棘轮}。模块化、标准化和淘汰旧接口也能让复杂度净下降，具体方向要从架构与运行记录判断。

这个过程可以写成一个递推关系
\begin{equation}
C_{n+1}=(1+\alpha)\,C_n+c
\label{eq:ratchet-recursion}
\end{equation}
其中$C_n$为按一致口径统计的复杂度，$c$为新增层的直接增量，$\alpha$为既有结构因新层接入而增加的比例。固定$c>0$、$\alpha>0$时，这个递推给出指数增长；$\alpha=0$时为线性增长。它描述依赖持续累积的一种情形，不能由层数本身推定实际系统服从固定比例。

把 \eqref{eq:ratchet-recursion} 展开，可以看清棘轮的形状。设起始复杂度为 $C_{0}$，$n$ 层之后
\begin{equation}
C_{n}=(1+\alpha)^{n}C_{0}+\frac{c}{\alpha}\bigl[(1+\alpha)^{n}-1\bigr],
\qquad
C^{\min}_{n}=\beta\,C_{n},\qquad 0<\beta<1
\label{eq:ratchet-floor}
\end{equation}
$\alpha>0$时第一式给出递推解；$\alpha=0$时取极限$C_n=C_0+nc$。第二式把拆除后仍保留的复杂度记为当前值的$\beta$倍，$\beta$需要由具体退出方案估计。旧依赖可以被移交、重构或清除，因此$C_n^{\min}$是情景下的可达水平，不是普遍物理下限。

这个差就是锁定带，它的累积量就是卸载前 $n$ 层要付的兼容税
\begin{equation}
K_{n}=\sum_{k=1}^{n}\bigl(C_{k}-C^{\min}_{k}\bigr)=(1-\beta)\sum_{k=1}^{n}C_{k}
\label{eq:compat-tax}
\end{equation}
$K_n$把逐层差额汇总成一个兼容负担指标，单位随复杂度口径确定。各层有共享接口时，直接相加会重复计算；转换为退出费用，需要分别估计拆除、迁移、验证与替代支持。实际成本也可能在集中重构时下降，不能只凭这个和式断言退出账单必然指数上升。

自动驾驶正是覆盖层困境的典型案例\cite{kalra2016}。道路原本不是为机器视觉设计的，交通规则原本不是为算法责任设计的，保险制度原本不是为无人主体设计的，行人和司机也不会完全按照机器协议行动。自动驾驶如果要真正简单，最合理的办法是重构环境：封闭道路、人车分流、统一车路通信、标准化交通参与者、清晰责任边界，但这意味着重构城市交通的底层，代价极高。于是自动驾驶选择覆盖层路线：不改造所有道路，而是让车携带更多传感器；不要求所有人类行为标准化，而是让算法处理长尾；不立即重写责任制度，而是叠加安全员、远程接管、监管沙盒、保险补丁和数据记录。环境没有变简单，车辆只能变复杂；道路没有被低熵化，算法只能承受更多不确定性。自动驾驶看似是机器替代司机，深层上却是把道路重构成本转移到了单车智能、感知冗余、云端训练和责任系统中。这也是为什么人机混行阶段格外困难：纯人类驾驶是一种相对稳定的旧秩序，纯机器驾驶在高度标准化环境中也可能稳定，最复杂的是二者混合的过渡态，人类、机器、旧道路、旧法规、非标准行为、保险责任和社会预期同时存在，系统进入高熵混合区。此时智能越强，不一定系统越简单，很多算力其实是在为旧环境的无序性买单。
\begin{figure}[ht]
\centering
% >>>>>>>>>> figs/fig-ratchet.tex >>>>>>>>>>
% 图 3-3  复杂度棘轮：能加不能减的那一部分
% 横轴时间，纵轴系统复杂度。实线是当前复杂度，每加一层上一个台阶；
% 虚线是可回退的下限——它也在涨，只是涨得慢。两线之间是锁定带：
% 想退回上一代，得先付清这一带里的兼容税。台阶越多，带越宽。
% 末端点划线是「只覆盖不重构」的外推：层继续堆，可回退下限抬得更快。
\begin{tikzpicture}[x=1cm, y=1cm, font=\small,
  ax/.style={-{Stealth[length=2.0mm]}, ink, line width=0.6pt},
  arr/.style={{Stealth[length=1.5mm]}-{Stealth[length=1.5mm]}, ink!60, line width=0.5pt}
]
\useasboundingbox (-1.05,0.80) rectangle (12.85,-6.05);

\def\aX{1.15} \def\aY{-4.45} \def\dx{0.98} \def\dy{0.40} \def\aR{11.30}
% 坐标宏里不写括号，见 fig-three-media.tex 的注释
\def\px#1{\aX+\dx*#1}
\def\py#1{\aY+\dy*#1}

% ---------- 锁定带：上界走台阶，下界走缓坡，两界之间填色 ----------
\fill[rust!13]
  (\px{0.4},\py{0.55}) -- (\px{1},\py{0.55}) -- (\px{1},\py{1.90})
  -- (\px{2},\py{1.90}) -- (\px{2},\py{3.00}) -- (\px{3},\py{3.00})
  -- (\px{3},\py{4.10}) -- (\px{4},\py{4.10}) -- (\px{4},\py{5.30})
  -- (\px{5},\py{5.30}) -- (\px{5},\py{6.40}) -- (\px{6},\py{6.40})
  -- (\px{6},\py{7.60}) -- (\px{7},\py{7.60}) -- (\px{7},\py{8.80})
  -- (\px{8},\py{8.80}) -- (\px{8},\py{10.10}) -- (\px{9.6},\py{10.10})
  -- (\px{9.6},\py{1.62})
  -- (\px{8},\py{1.42}) -- (\px{6},\py{1.18}) -- (\px{4},\py{0.95})
  -- (\px{2},\py{0.85}) -- (\px{0.4},\py{0.40}) -- cycle;

% ---------- 坐标轴 ----------
\draw[ax] (\aX,\aY) -- (\aR,\aY) node[anchor=west, font=\scriptsize, text=ink] {时间};
\draw[ax] (\aX,\aY) -- (\aX,\py{10.9})
      node[anchor=south, font=\scriptsize, text=ink, yshift=1pt] {系统复杂度};

% ---------- 实线：当前复杂度 ----------
\draw[deepblue!85, line width=1.9pt]
  (\px{0.4},\py{0.55}) -- (\px{1},\py{0.55}) -- (\px{1},\py{1.90})
  -- (\px{2},\py{1.90}) -- (\px{2},\py{3.00}) -- (\px{3},\py{3.00})
  -- (\px{3},\py{4.10}) -- (\px{4},\py{4.10}) -- (\px{4},\py{5.30})
  -- (\px{5},\py{5.30}) -- (\px{5},\py{6.40}) -- (\px{6},\py{6.40})
  -- (\px{6},\py{7.60}) -- (\px{7},\py{7.60}) -- (\px{7},\py{8.80})
  -- (\px{8},\py{8.80}) -- (\px{8},\py{10.10}) -- (\px{9.6},\py{10.10});
% ---------- 虚线：可回退下限 ----------
\draw[forest!85, line width=1.4pt, dashed]
  (\px{0.4},\py{0.40}) -- (\px{2},\py{0.85}) -- (\px{4},\py{0.95})
  -- (\px{6},\py{1.18}) -- (\px{8},\py{1.42}) -- (\px{9.6},\py{1.62});

% ---------- 各层之间的分界（细点线，只在锁定带里画）----------
\foreach \k in {1,3,5,7} {
  \draw[slate!45, line width=0.4pt, densely dotted]
        (\px{\k},\py{0.8}) -- (\px{\k},\py{9.2});
}

% ---------- 注记 ----------
\node[font=\scriptsize, text=deepblue!85, anchor=south west] at (\px{4.35},\py{5.42})
      {当前复杂度：每加一层，上一个台阶};
\node[font=\scriptsize, text=forest!85, anchor=north west] at (\px{0.45},\py{0.30})
      {可回退下限};
\node[font=\scriptsize, text=slate, anchor=south] at (\px{1.5},\py{1.98}) {第一层};
\node[font=\scriptsize, text=slate, anchor=south] at (\px{4.5},\py{4.18}) {第二层};
\node[font=\scriptsize, text=slate, anchor=south] at (\px{6.5},\py{6.48}) {第三层};
\node[font=\scriptsize, text=slate, anchor=south] at (\px{7.6},\py{8.88}) {第四层};

% 锁定带厚度
\draw[arr] (\px{9.30},\py{1.62}) -- (\px{9.30},\py{10.10});
\node[font=\scriptsize, text=ink, anchor=west, align=left] at (\px{9.40},\py{6.40})
      {锁定带\\（兼容税）};

% 末端外推
\draw[ink!60, line width=1.0pt, densely dashdotted]
      (\px{9.6},\py{10.10}) -- (\px{10.7},\py{11.35});
\node[font=\scriptsize, text=ink, anchor=north west, align=left]
      at (\px{9.65},\py{10.02}) {只覆盖\\不重构};
\end{tikzpicture}

% <<<<<<<<<< figs/fig-ratchet.tex <<<<<<<<<<

\caption{复杂度棘轮。两条线的间距即锁定带：想退回上一代，得先付清这一带里的兼容税。带的两侧都在上涨，但涨速不同，间距因此逐级拉开；台阶越多，带越宽。}
\label{fig:ratchet}
\end{figure}


\subhead{兼容税与落地税}

具身智能面对的是同一种赛博寄生，只是空间更广。家庭、医院、餐厅、仓库、商场和街道，都不是为机器人身体设计的：门把手默认人手，楼梯默认双腿，厨房默认眼手协调，工具默认人的触觉和常识，家庭物品默认人类的模糊处理能力。机器人要进入这些空间，就必须支付兼容税：更多自由度、更复杂传感器、更高实时控制、更强安全冗余、更昂贵维护体系。工业机器人之所以成功，是因为它们通常被放在低熵环境中，固定位置、重复动作、标准物料、隔离空间、明确任务；通用具身智能之所以困难，是因为它要进入高熵人类环境，这里物体位置不固定，任务目标含糊，空间布局多样，人类行为不可预测，接触反馈复杂。给机器人装上大模型，并不会自动取消摩擦、碰撞、误差、遮挡、材料差异和责任风险。

这种成本分布有一个更早的表述。Hans Moravec 在 1988 年指出，让计算机在智力测验或下棋上达到成人水平相对容易，要让它在感知和运动上达到一岁儿童的水平却极其困难\cite{moravec1988}。这个判断在三十年后的具身智能研究中仍然成立：需要显式推理的任务可以被形式化、被标注、被搜索，而抓取、行走、避障这类能力依赖大量隐性调节，只能在接触中训练出来，训练成本恰好落在最不容易复制的那一部分。兼容税之所以高昂，不是因为机器人不够聪明，而是因为它要补的正是人类身体用漫长演化预先支付掉的那部分调节。

所谓“落地税”，就发生在模型接触现实的一瞬间。任何世界模型都是现实的压缩，而现实总有残差，模拟器中的训练无法穷尽真实世界的高频噪声、边缘情况、材料变化和人类偶然性。机器人一旦行动，就会改变被建模对象：拿起杯子，杯子位置变了；推开门，门的状态变了；进入房间，人的行为也可能变了。物理AI不是在静态地图上执行，而是在被自己扰动的环境中连续修正。因此，覆盖层物理AI的真正难点，不是智能不足，而是旧世界没有为智能低成本行动而重构。AI必须在旧世界中支付兼容税、残差税和维护税。它通过比特层提高局部效率，却不得不承担原子层的沉重、混乱和不可回滚后果。

\subhead{覆盖与重构之间}

覆盖层策略的问题，不在于它无效，而在于它容易被误认为重构。覆盖可以延长旧系统生命，可以提高局部效率，可以让旧基础设施获得新的控制能力，但覆盖不能无限叠加。每加一层，都要问：它减少的摩擦，是否大于它增加的复杂度？它节省的人力，是否小于它带来的维护负担？它提升的效率，是否值得它增加的脆弱性？一个系统如果只覆盖、不重构，就会逐渐形成高效而脆弱的复杂网络。正常运行时，它显得极其智能；一旦接口失效、数据异常、传感器故障、平台宕机或责任链断裂，复杂性就会反过来成为风险放大器。旧系统的问题没有消失，只是被更多智能层遮住；旧责任没有消失，只是被平台、算法、运营者、用户、监管者和人工兜底重新分散。

所以，物理AI的成熟，不只是让模型更大、传感器更多、控制更精细，而是判断何时覆盖，何时重构，何时放弃。低风险、低耦合、高标准化场景适合覆盖：封闭园区、工业产线、仓储物流、固定路线运输，都可以通过覆盖层获得明显收益。高风险、高开放、高耦合场景不能只靠覆盖：自动驾驶、医疗护理、家庭机器人、城市治理，最终都需要环境设计、制度重构和责任边界的同步变化。未来的物理AI竞争，不只是模型能力之争，而是环境低熵化能力、接口标准能力、基础设施重构能力和责任制度设计能力之争：谁能减少兼容税，谁就能让智能行动更便宜；谁能降低环境混乱度，谁就能让具身系统更可靠；谁能识别覆盖与重构的边界，谁就能避免复杂度棘轮把系统拖向脆弱。赛博寄生是物理AI的起点，不是终点。AI进入原子世界，往往先作为覆盖层寄生在旧系统之上；但如果永远停留在覆盖层，文明只会得到越来越复杂的补丁网络。真正可靠的物理智能，不是无止境地给旧世界贴上传感器和算法，而是在覆盖、重构和保留之间找到最低总黏滞的路径。

旧世界不会因为智能到来而自动消失。它会先被覆盖、被监控、被调度、被补丁化。物理AI的第一性现实就在这里：智能要改变原子世界，必须先寄生原子世界；而它最终能否成熟，取决于它是否有能力从寄生走向真正的环境重构。但无论选择覆盖还是重构，智能都在持续向一个高黏滞系统输入。系统未必按输入的方向改变：它可能沿阻力更低的路径绕开，可能把吸收不了的输入堆在自己内部，也可能因为反馈过急而开始摆动。这三种反应不受策略选择支配，只受系统自身动力学的支配。
